% ===== PRÉAMBULE CANONIQUE — à coller VERBATIM en tête de CHAQUE .tex =====
\documentclass[11pt,a4paper]{article}
\usepackage[utf8]{inputenc}\usepackage[T1]{fontenc}\usepackage{lmodern}
\usepackage[french]{babel}
\usepackage{amsmath,amssymb,mathtools}\usepackage{icomma}
\usepackage{siunitx}\sisetup{locale=FR}  % \qty{}{}, \si{}, \ang{}
\usepackage{xcolor}\usepackage{colortbl}
\usepackage{tikz}\usepackage{pgfplots}\pgfplotsset{compat=1.18}
\usepackage[siunitx,europeanresistors]{circuitikz} % circuits (élec.)
\usepackage[most]{tcolorbox}\usepackage{enumitem}\usepackage{array,booktabs}
\usepackage[a4paper,margin=2.2cm]{geometry}
\usetikzlibrary{arrows.meta,calc,angles,quotes,positioning,patterns,%
  backgrounds,shapes.geometric,decorations.pathmorphing,decorations.markings,intersections}
\definecolor{primary}{RGB}{222,49,99}\definecolor{primarydark}{RGB}{150,22,55}
\definecolor{defcol}{RGB}{0,114,178}\definecolor{methcol}{RGB}{52,92,148}
\definecolor{excol}{RGB}{0,135,90}\definecolor{attncol}{RGB}{200,40,55}
\tcbset{boxbase/.style={enhanced,breakable,boxrule=0.9pt,arc=2.5pt,
  left=3mm,right=3mm,top=2.3mm,bottom=2.3mm,fonttitle=\bfseries,coltitle=white}}
% --- boîtes TUTORIEL (noms EXACTS, ne pas renommer) ---
\newtcolorbox{cle}[1][]{boxbase,colback=defcol!6,colframe=defcol,
  title={\textsc{À retenir}\ifx\relax#1\relax\relax\else\textmd{ : \itshape #1}\fi}}
\newtcolorbox{methode}[1][]{boxbase,colback=methcol!7,colframe=methcol,sharp corners,
  title={\textsc{Méthode}\ifx\relax#1\relax\relax\else\textmd{ --- \itshape #1}\fi}}
\newtcolorbox{exemplebox}[1][]{boxbase,colback=excol!5,colframe=excol,
  title={\textsc{Exemple}\ifx\relax#1\relax\relax\else\textmd{ : \itshape #1}\fi}}
\newtcolorbox{piege}[1][]{boxbase,colback=attncol!6,colframe=attncol,
  title={\textsc{Piège}\ifx\relax#1\relax\relax\else\textmd{ : \itshape #1}\fi}}
\newtcolorbox{remarque}[1][]{boxbase,colback=black!4,colframe=black!45,
  title={\textsc{Remarque}\ifx\relax#1\relax\relax\else\textmd{ : \itshape #1}\fi}}
% --- boîtes EXERCICE / CORRIGÉ (noms EXACTS, auto-comptées) ---
\newtcolorbox[auto counter]{exo}[1][]{boxbase,colback=primary!4,colframe=primary,
  title={\textsc{Exercice}~\thetcbcounter\ifx\relax#1\relax\relax\else\textmd{ --- \itshape #1}\fi}}
\newtcolorbox[auto counter]{corrige}[1][]{boxbase,colback=excol!4,colframe=excol,
  title={\textsc{Corrigé}~\thetcbcounter\ifx\relax#1\relax\relax\else\textmd{ --- \itshape #1}\fi}}
\newcommand{\vv}[1]{\vec{#1}}\newcommand{\ds}{\displaystyle}
\newcommand{\R}{\mathbb{R}}\newcommand{\N}{\mathbb{N}}\newcommand{\Z}{\mathbb{Z}}
% (un \usepackage{fancyhdr}+en-tête est possible ; il est ignoré côté web)
% =========================================================================
\begin{document}
\sisetup{per-mode=symbol}

\begin{exo}[de l'électron à la charge]
La charge élémentaire vaut $e = \num{1.6e-19}\ \si{\coulomb}$.
\begin{enumerate}[label=\alph*)]
  \item Un corps a \emph{perdu} $N = \num{5e12}$ électrons. Quelle est sa charge $q$ (valeur et
        signe) ?
  \item Un autre corps a \emph{gagné} $N = \num{2.5e13}$ électrons. Quelle est sa charge $q$ ?
\end{enumerate}
\end{exo}

\begin{exo}[compter les électrons en excès]
Une petite sphère porte une charge $q = \SI{-8e-9}{\coulomb}$.
\begin{enumerate}[label=\alph*)]
  \item Combien d'électrons en excès (ou en défaut) porte-t-elle ?
  \item S'agit-il d'un excès ou d'un défaut d'électrons ? Justifie par le signe de $q$.
\end{enumerate}
\end{exo}

\begin{exo}[conducteur ou isolant ?]
On dispose des matériaux suivants : \emph{cuivre, verre, plastique, argent, ébonite, aluminium}.
\begin{enumerate}[label=\alph*)]
  \item Range-les en deux colonnes : \textbf{conducteurs} / \textbf{isolants}.
  \item Lesquels peut-on électriser \emph{durablement par frottement} ? Explique en une phrase.
\end{enumerate}
\end{exo}

\begin{exo}[électrisation par contact]
Une sphère métallique \textbf{chargée positivement} (tenue par un manche isolant) touche une sphère
métallique \textbf{neutre}, identique, puis on les sépare.
\begin{center}
\begin{tikzpicture}[scale=0.9,>=Stealth,
    bille/.style={circle,draw=black!55,line width=0.8pt,minimum size=13mm,inner sep=0pt,font=\small}]
  \node[bille,fill=attncol!18] (A) at (0,0) {$+$};
  \node[font=\scriptsize,below=1pt] at (A) {chargée $+$};
  \node[bille,fill=black!5] (B) at (2,0) {$0$};
  \node[font=\scriptsize,below=1pt] at (B) {neutre};
  \draw[->,thick] (3.3,0) -- (4.4,0) node[midway,above,font=\scriptsize]{contact};
  \node[bille,fill=black!8] (C) at (5.8,0) {?};
  \node[bille,fill=black!8] (D) at (7.6,0) {?};
  \node[font=\scriptsize,below=1pt] at ($(C)!0.5!(D)+(0,-0.55)$) {après séparation};
\end{tikzpicture}
\end{center}
\begin{enumerate}[label=\alph*)]
  \item Quel est le signe de chaque sphère après le contact ?
  \item Ce mode d'électrisation crée-t-il de la charge ? Quel principe l'interdit ?
\end{enumerate}
\end{exo}

\begin{exo}[influence : le conducteur reste neutre]
On approche (sans le toucher) un bâton \textbf{chargé négativement} d'une barre métallique neutre.
\begin{center}
\begin{tikzpicture}[scale=0.9,>=Stealth,
    pe/.style={circle,draw=black!50,minimum size=4.5mm,inner sep=0pt,font=\scriptsize,text=white}]
  \draw[fill=defcol!22,line width=0.9pt] (-3.6,-0.35) rectangle (-2.2,0.35);
  \foreach \x in {-3.3,-2.95,-2.6}{\node[pe,fill=defcol] at (\x,0) {$-$};}
  \node[defcol,font=\scriptsize,above] at (-2.9,0.45) {bâton $(-)$};
  \draw[->,defcol,dashed] (-2.15,0) -- (-1.15,0);
  \draw[fill=black!6,line width=0.9pt] (-1,-0.45) rectangle (2.4,0.45);
  \node[font=\scriptsize] at (0.7,-0.75) {barre métallique neutre};
\end{tikzpicture}
\end{center}
\begin{enumerate}[label=\alph*)]
  \item De quel côté de la barre s'accumulent les charges $+$ ? les charges $-$ ? (complète
        mentalement le schéma)
  \item La barre est-elle globalement chargée ? Justifie.
\end{enumerate}
\end{exo}

\end{document}
